\documentclass{beamer}

\usetheme{Berlin}
\usepackage{polski}
\usepackage{times}
\usepackage[T1]{fontenc}
\usepackage{graphicx}
\usepackage{listings}
\usepackage{hyperref}
\usepackage{subcaption}
\usefonttheme{serif}

\title{Egzamin dyplomowy}
\subtitle{Układ aktywnego tłumienia hałasu pod nadzorem jądra systemu Linux z~rozszerzeniami real-time}

\author{Daniel Kopiecki}
\institute{Politechnika Śląska \\ Wydział Automatyki, Elektroniki i Informatyki}

\date{\today}

\begin{document}

\begin{frame}
    \titlepage
\end{frame}

\begin{frame}
    \frametitle{Cel pracy}
    Skonstruowanie prostego układu aktywnej redukcji hałasu w oparciu o
    samodzielnie zaprojektowaną elektronikę, wykorzystując system operacyjny
    ogólnego przeznaczenia z jądrem zawierającym rozszerzenia do pracy w czasie
    rzeczywistym. 
    \vfill
    \par Projekt układu można podzielić na:
    \begin{enumerate}
        \item \textbf{Część sprzętową} --- płytkę drukowaną pracującą jako
        zewnętrzna karta pomiarowa.
        \item \textbf{Część programową} --- zmodyfikowany system operacyjny wraz
        z aplikacją działającą w przestrzeni użytkownika.
    \end{enumerate}
\end{frame}

\begin{frame}
    \frametitle{Najważniejsze założenia projektowe}
    \begin{itemize}
        \item Wykonanie analogowego układu elektronicznego w postaci płytki
        drukowanej, zdolnego do przetworzenia dwóch wejściowych i dwóch
        wyjściowych kanałów audio do postaci cyfrowej wraz z ich dostarczeniem
        do wybranej jednostki obliczeniowej.
        \item Wykorzystanie jądra wyposażonego w rozszerzenia do pracy w czasie
        rzeczywistym wraz z weryfikacją ich działania.
        \item Implementacja aplikacji w przestrzeni użytkownika realizującej
        proces aktywnej redukcji hałasu wraz z narzędziami do testów i
        symulacji.
        \item Wykorzystanie ogólnodostępnych i otwartoźródłowych rozwiązań.
    \end{itemize}
\end{frame}

\begin{frame}
    \frametitle{Architektura systemu}
    \begin{figure}[H]
        \centering
        \includegraphics[width=0.5\linewidth]{../images/top.png}
        \caption{Architektura skonstruowanego systemu}
    \end{figure}
\end{frame}

\begin{frame}
    \frametitle{Część sprzętowa}
    \begin{figure}[H]
        \centering
        \includegraphics[width=\linewidth]{../images/hw_top.png}
        \caption{Część sprzętowa --- architektura}
    \end{figure}
\end{frame}

\begin{frame}
    \frametitle{Część sprzętowa --- zasilania}
    \begin{figure}[H]
        \centering
        \includegraphics[width=0.75\linewidth]{../images/power_dist.png}
        \caption{Diagram zasilań części sprzętowej}
    \end{figure}
\end{frame}

\begin{frame}
    \frametitle{Część sprzętowa --- tory analogowe}
    \begin{figure}[H]
        \centering
        \includegraphics[width=\linewidth]{../images/analog_in.png}
        \caption{Wejściowe tory analogowe}
    \end{figure}
    \begin{figure}[H]
        \centering
        \includegraphics[width=\linewidth]{../images/analog_out.png}
        \caption{Wyjściowe tory analogowe}
    \end{figure}
\end{frame}

\begin{frame}
    \frametitle{Część sprzętowa --- przetworniki}
    \begin{figure}[H]
        \centering
        \includegraphics[width=\linewidth]{../images/adc_dac.png}
        \caption{Sekcja przetworników}
    \end{figure}
\end{frame}

\begin{frame}
    \frametitle{Część programowa}
    \begin{figure}[H]
        \centering
        \includegraphics[width=0.6\linewidth]{../images/app_top.png}
        \caption{Architektura programowa}
    \end{figure}
\end{frame}

\begin{frame}
    \frametitle{Część programowa --- aplikacja}
    \begin{figure}[H]
        \centering
        \includegraphics[width=\linewidth]{../images/app_flow.png}
        \caption{Przebieg działania aplikacji}
    \end{figure}
\end{frame}

\begin{frame}
    \frametitle{Rezultaty --- karta pomiarowa}
    \begin{figure}[H]
        \centering
        \includegraphics[width=0.8\linewidth]{../images/board.png}
        \caption{Karta pomiarowa}
    \end{figure}
\end{frame}

\begin{frame}

    \frametitle{Rezultaty --- pomiary z wykorzystaniem karty}
    \centering
    \begin{minipage}{.49\textwidth}
    \centering
    \begin{figure}
        \begin{subfigure}{0.75\textwidth}
          \centering
          \includegraphics[width=\linewidth]{../images/200hz_in.png}
        \end{subfigure}
        \begin{subfigure}{\textwidth}
          \centering
          \includegraphics[width=\linewidth]{../images/200hz_psd_in.png}
        \end{subfigure}
        \caption{Sinusoida 200 Hz}
        \end{figure}
    \end{minipage}%
    \begin{minipage}{.49\textwidth}
        \begin{figure}
            \begin{subfigure}{0.75\textwidth}
              \centering
              \includegraphics[width=\linewidth]{../images/500hz_in.png}
            \end{subfigure}
            \begin{subfigure}{\textwidth}
              \centering
              \includegraphics[width=\linewidth]{../images/500hz_psd_in.png}
            \end{subfigure}
            \caption{Sinusoida 500 Hz}
            \end{figure}
    \end{minipage}
    
\end{frame}

\begin{frame}
\frametitle{Stanowisko pomiarowe}
\begin{figure}[H]
    \centering
    \includegraphics[width=0.8\linewidth]{../images/stand.png}
    \caption{Stanowisko pomiarowe}
\end{figure}
\end{frame}

\begin{frame}
\frametitle{Przykładowe przebiegi błedu dla sygnałów testowych}
\begin{figure}[H]
    \centering
    \begin{subfigure}[b]{0.49\textwidth}
        \centering
        \includegraphics[width=\textwidth]{../images/alpha_120hz_01.png}
    \end{subfigure}
    \begin{subfigure}[b]{0.49\textwidth}
        \centering
        \includegraphics[width=\textwidth]{../images/alpha_220hz_01.png}
    \end{subfigure}
    \begin{subfigure}[b]{0.49\textwidth}
        \centering
        \includegraphics[width=\textwidth]{../images/alpha_sum_01.png}
    \end{subfigure}
       \caption{Przebiegi referencyjne $e(n)$ dla sygnałów testowych.}
\end{figure}
\end{frame}

\begin{frame}
    \frametitle{Przebiegi błędu --- jądro bez RT}
    \begin{figure}[H]
        \centering
        \begin{subfigure}[b]{0.49\textwidth}
            \centering
            \includegraphics[width=\textwidth]{../images/bad_120.png}
        \end{subfigure}
        \begin{subfigure}[b]{0.49\textwidth}
            \centering
            \includegraphics[width=\textwidth]{../images/bad_220.png}
        \end{subfigure}
        \begin{subfigure}[b]{0.49\textwidth}
            \centering
            \includegraphics[width=\textwidth]{../images/bad_sum.png}
        \end{subfigure}
        \caption{Przebiegi referencyjne $e(n)$ dla sygnałów testowych.}
    \end{figure}
\end{frame}

\begin{frame}
    \frametitle{Losowe pobudzenia sinusoidalne}
    \begin{figure}[H]
        \centering
        \includegraphics[width=0.9\linewidth]{../images/anc_rand_sin.png}
        \caption{Eksperyment --- losowe pobudzenia sinusoidalne}
    \end{figure}
\end{frame}

\begin{frame}
    \frametitle{Hałas szerokopasmowy}
    \begin{figure}[H]
        \centering
        \includegraphics[width=0.9\linewidth]{../images/anc_rand_0.png}
        \caption{Eksperyment --- hałas losowy}
    \end{figure}
\end{frame}

\begin{frame}
    \frametitle{Hałasy rzeczywiste --- młyny kulowe}
    \begin{figure}[H]
        \centering
        \includegraphics[width=0.9\linewidth]{../images/mlyny.png}
        \caption{Eksperyment --- młyny kulowe}
    \end{figure}
\end{frame}

\begin{frame}
    \frametitle{Hałasy rzeczywiste --- pompy wody}
    \begin{figure}[H]
        \centering
        \includegraphics[width=0.9\linewidth]{../images/anc_opc.png}
        \caption{Eksperyment --- pompy wody}
    \end{figure}
\end{frame}

\begin{frame}
    \frametitle{Badania}
    \begin{itemize}
        \item \textbf{Parametry} --- wielkość kroku adaptacji, długość filtra,
        częstotliwość próbkowania, długość fragmentu, poziom priorytetu,
        wytyczna planisty.
        \item \textbf{Rezultaty} --- system jest w stanie tłumić sygnały
        wąskopasmowe o około 15 dB do założonego pasma (500 Hz), jednak nie
        radzi sobie z szerokopasmowymi.
        \item \textbf{Powody} --- duże opóźnienia, buforowanie, zbyt wolna
        adaptacja, latencja systemu operacyjnego.
        \item \textbf{Możliwe usprawnienia} --- uproszczenie elektroniki,
        nadpróbkowanie, leaky LMS, urównoleglenie niektórych operacji.
    \end{itemize}
\end{frame}

\begin{frame}
    \huge{Serdecznie dziękuję za uwagę}
\end{frame}

\end{document}

